% Generated draft: requires mathematical and source review.
\documentclass[11pt]{article}
\usepackage[margin=25mm]{geometry}
\usepackage{fontspec}
\setmainfont{Latin Modern Roman}
\usepackage{amsmath,amssymb,mathtools,unicode-math}
\setmathfont{Latin Modern Math}
\usepackage{enumitem,needspace,xurl,hyperref}
\setlength{\parindent}{0pt}
\setlength{\parskip}{5pt}
\setlength{\emergencystretch}{3em}
\newcommand{\N}{\mathbb{N}}
\newcommand{\Z}{\mathbb{Z}}
\newcommand{\Q}{\mathbb{Q}}
\newcommand{\R}{\mathbb{R}}
\newcommand{\C}{\mathbb{C}}
\newcommand{\F}{\mathbb{F}}
\newcommand{\A}{\mathbb{A}}
\newcommand{\PP}{\mathbb{P}}
\newcommand{\E}{\mathbb{E}}
\newcommand{\eps}{\varepsilon}
\newcommand{\dd}{\,\mathrm{d}}
\begin{document}
\section*{UnsolvedMath ID: 151}
\textbf{Source problem code:} catalogue entry
\subsection{Definitions and mathematical statement}
The source asks for conditions on a set $A$ under which the number of solutions of
\[F(x_1,\ldots,x_n)=C\]
with all coordinates in $A$ is roughly $\alpha^n$ times the corresponding number with all coordinates in $[X]$. For finite sets $S$ and $A\subseteq S$, and an equation defined on $S^n$, the relevant notation would be
\[N_{F,C}(A)=\bigl|\{(x_1,\ldots,x_n)\in A^n:F(x_1,\ldots,x_n)=C\}\bigr|,\qquad \alpha=\frac{|A|}{|S|}.\]
Thus the requested comparison has the schematic form
\[N_{F,C}(A)\approx \alpha^n N_{F,C}(S),\]
with $S$ represented by $[X]$ in the imported wording. This notation only explains the counting question; it does not fix its missing hypotheses. In particular, the record does not specify the ambient number system and box, the precise meaning of $[X]$, the class and coefficients of $F$, the allowed constants $C$, the range of $n$, any nonsingularity or local-solubility hypotheses, whether repeated coordinates are allowed, or the density and uniformity conditions on $A$. It also gives no limiting process or error bound defining $\approx$ (``roughly''). A finite-ambient density definition above is conditional notation, not a claim that this choice exhausts the intended setting.

The supplied triage checked 2026-08-17 states that circle-method, transference, and pseudorandomness results address many precise variants, but that the imported wording does not determine one literature-checkable theorem. It assigns partially solved / PARTIAL-PROGRESS to the record; the separate registry field is open. No specific variant or new theorem-level status is asserted here.
\subsection{Short English statement}
Find hypotheses making a set of density $\alpha$ contain about the random-model proportion $\alpha^n$ of the solutions to $F(x_1,\ldots,x_n)=C$. The ambient domain, polynomial hypotheses, limiting parameter, and permitted error must be fixed before this becomes a precise mathematical problem.
\subsection{Sources}
\begin{itemize}
\item Local UnsolvedMath registry, permanent ID 151, source code catalogue entry, ``Solutions to Polynomial Equations in Dense Sets,'' in Ben Green's 100 Open Problems; supplied statement and background.
\item B. Green, \emph{100 Open Problems}, maintained author notes accessed 2026-08-17, supplied authoritative secondary reference: \url{https://people.maths.ox.ac.uk/greenbj/papers/open-problems.pdf}. The source record does not identify a specific numbered passage here.
\item Supplied literature triage checked 2026-08-17, PARTIAL-PROGRESS. It explicitly retains missing formal data and reports no OCR or statement repair; dated triage does not exclude later or unindexed results.
\end{itemize}
\end{document}
